
\documentclass[10pt,oneside]{book}

% ==================================================
% PORQUERÍA MATEMÁTICA
% Guía de supervivencia matemática
% ==================================================

\usepackage[utf8]{inputenc}
\usepackage[T1]{fontenc}
\usepackage[spanish,mexico]{babel}
\usepackage{lmodern}
\usepackage{microtype}
%\usepackage[a4paper,margin=1.65cm,headheight=14pt]{geometry}
\usepackage[
paperwidth=5.5in,
paperheight=8.5in,
top=0.65in,
bottom=0.65in,
inner=0.75in,
outer=0.55in
]{geometry}
\usepackage{amsmath,amssymb,mathtools}
\usepackage{graphicx}
\usepackage{booktabs,array,tabularx,multicol}
\usepackage[most]{tcolorbox}
\usepackage{xcolor}
\usepackage{enumitem}
\usepackage{fancyhdr}
\usepackage{hyperref}
\usepackage{titlesec}

\definecolor{tinta}{HTML}{202020}
\definecolor{gris}{HTML}{EEEEEE}
\definecolor{rojo}{HTML}{8D2424}
\definecolor{azul}{HTML}{244A68}

\hypersetup{
	colorlinks=true,
	linkcolor=azul,
	urlcolor=rojo,
	pdftitle={Porquería Matemática},
	pdfauthor={Brillinury y Samuel}
}

\pagestyle{fancy}
\fancyhf{}
\fancyhead[L]{\small\textit{Porquería Matemática}}
\fancyhead[R]{\small Guía de supervivencia}
\fancyfoot[C]{\thepage}
\setlength{\parindent}{0pt}
\setlength{\parskip}{3pt}
\setlist{nosep,leftmargin=*}
\setcounter{tocdepth}{1}

\titleformat{\chapter}
{\Large\bfseries\color{tinta}}
{\thechapter.}{0.5em}{}
\titleformat{\section}
{\large\bfseries\color{azul}}
{\thesection}{0.5em}{}

\newtcolorbox{formula}{
	colback=gris,colframe=tinta,
	boxrule=.5pt,arc=1mm,
	left=2mm,right=2mm,top=1mm,bottom=1mm
}
\newtcolorbox{ejemplo}{
	colback=white,colframe=azul,
	boxrule=.6pt,arc=1mm,
	title=Ejemplo resuelto,
	fonttitle=\bfseries
}
\newtcolorbox{ojo}{
	colback=red!4,colframe=rojo,
	boxrule=.5pt,arc=1mm,
	title=No lo hagas peor,
	fonttitle=\bfseries
}

\newcommand{\R}{\mathbb{R}}
\newcommand{\N}{\mathbb{N}}
\newcommand{\Z}{\mathbb{Z}}
\newcommand{\sen}{\operatorname{sen}}
\newcommand{\Prob}{\operatorname{P}}

\begin{document}
	
	% ==================================================
	% PORTADA
	% ==================================================
	
	\begin{titlepage}
		\centering
		\vspace*{2cm}
		
		{\Huge\bfseries PORQUERÍA\\[3mm] MATEMÁTICA\par}
		
		\vspace{1cm}
		{\Large Guía de supervivencia escolar\par}
		
		\vspace{1.5cm}
		\begin{tcolorbox}[
			width=.85\textwidth,
			colback=gris,colframe=tinta,
			arc=2mm,boxrule=1pt
			]
			\centering
			Una guía rápida de matemáticas para sobrevivir
			a la primaria, secundaria y preparatoria.
			
			Sin revelaciones profundas. Sin discursos inspiradores.
			Con fórmulas, ejemplos y lo estrictamente necesario.
		\end{tcolorbox}
		
		\vfill
		
		{\large Autores: Brillinury y Samuel\par}
		
		\vspace{1cm}
		{\small Edición de consulta y regularización\par}
	\end{titlepage}
	
	\frontmatter
	\chapter*{Antes de empezar}
	\addcontentsline{toc}{chapter}{Antes de empezar}
	
	Este documento es un acordeón ampliado. No sustituye un curso
	completo ni garantiza que aparezcan todos los temas de tu examen.
	Los programas cambian según el grado, la escuela y el subsistema
	de bachillerato.
	
	La estrategia es sencilla:
	\begin{enumerate}
		\item Identifica qué tipo de problema tienes.
		\item Busca la fórmula correspondiente.
		\item Sustituye los datos.
		\item Haz las operaciones.
		\item Comprueba que el resultado tenga sentido.
	\end{enumerate}
	
	Las matemáticas no siempre son difíciles. A veces solamente
	están escritas de una forma poco amable.
	
	\tableofcontents
	
	\mainmatter
	
	% ==================================================
	\chapter{Aritmética: los números y sus problemas}
	
	\section{Las cuatro operaciones}
	
	\begin{itemize}
		\item \textbf{Suma:} juntar cantidades.
		\item \textbf{Resta:} calcular una diferencia.
		\item \textbf{Multiplicación:} sumar repetidamente una cantidad.
		\item \textbf{División:} repartir o calcular cuántas veces cabe algo.
	\end{itemize}
	
	\begin{formula}
		\[
		a+b,\qquad a-b,\qquad a\cdot b,\qquad \frac{a}{b}
		\]
		En una división, $b\neq 0$. Dividir entre cero no está definido.
	\end{formula}
	
	\begin{ejemplo}
		Calcula $18+4(3-1)$.
		
		Primero el paréntesis: $3-1=2$.
		
		Después la multiplicación: $4(2)=8$.
		
		Finalmente: $18+8=26$.
	\end{ejemplo}
	
	\section{Jerarquía de operaciones}
	
	El orden usual es:
	\begin{enumerate}
		\item Paréntesis y otros signos de agrupación.
		\item Potencias y raíces.
		\item Multiplicaciones y divisiones, de izquierda a derecha.
		\item Sumas y restas, de izquierda a derecha.
	\end{enumerate}
	
	\begin{ojo}
		No se resuelve necesariamente todo de izquierda a derecha.
		La jerarquía existe precisamente para evitar ese desastre.
	\end{ojo}
	
	\section{Divisibilidad y números primos}
	
	Un número es divisible entre otro si la división da un entero
	sin residuo.
	
	\begin{center}
		\begin{tabularx}{\textwidth}{cX}
			\toprule
			Divisor & Regla rápida\\
			\midrule
			2 & Termina en 0, 2, 4, 6 u 8.\\
			3 & La suma de sus cifras es múltiplo de 3.\\
			5 & Termina en 0 o 5.\\
			9 & La suma de sus cifras es múltiplo de 9.\\
			10 & Termina en 0.\\
			\bottomrule
		\end{tabularx}
	\end{center}
	
	Un número primo tiene exactamente dos divisores positivos:
	1 y él mismo. El 1 no es primo.
	
	\section{Mínimo común múltiplo y máximo común divisor}
	
	El \textbf{MCM} es el menor múltiplo positivo común.
	El \textbf{MCD} es el mayor divisor positivo común.
	
	Ejemplo: para 12 y 18,
	\[
	12=2^2\cdot3,\qquad 18=2\cdot3^2.
	\]
	\[
	\operatorname{MCD}(12,18)=2\cdot3=6
	\]
	\[
	\operatorname{MCM}(12,18)=2^2\cdot3^2=36.
	\]
	
	Para el MCD se toman los exponentes menores; para el MCM,
	los mayores.
	
	% ==================================================
	\chapter{Fracciones, decimales y porcentajes}
	
	\section{Fracciones}
	
	Una fracción representa partes de un todo:
	\[
	\frac{a}{b},\qquad b\neq0.
	\]
	$a$ es el numerador y $b$ el denominador.
	
	\begin{formula}
		\[
		\frac{a}{b}+\frac{c}{d}=\frac{ad+bc}{bd}
		\]
		\[
		\frac{a}{b}\cdot\frac{c}{d}=\frac{ac}{bd}
		\]
		\[
		\frac{a}{b}\div\frac{c}{d}=\frac{a}{b}\cdot\frac{d}{c}
		\]
		En la división también se requiere $c\neq0$.
	\end{formula}
	
	\begin{ejemplo}
		Calcula $\frac{1}{2}+\frac{1}{3}$.
		
		El denominador común puede ser 6:
		\[
		\frac{1}{2}+\frac{1}{3}
		=\frac{3}{6}+\frac{2}{6}
		=\frac{5}{6}.
		\]
	\end{ejemplo}
	
	\begin{ojo}
		No sumes los denominadores. En general,
		$\frac12+\frac13\neq\frac25$.
	\end{ojo}
	
	\section{Fracciones equivalentes y simplificación}
	
	Multiplicar o dividir numerador y denominador por el mismo
	número distinto de cero conserva el valor:
	\[
	\frac{2}{4}=\frac{1}{2}=\frac{5}{10}.
	\]
	
	Simplifica dividiendo arriba y abajo entre un divisor común.
	
	\section{Decimales}
	
	Para sumar y restar, alinea los puntos decimales.
	Para multiplicar, multiplica como si fueran enteros y luego
	coloca tantas cifras decimales como haya en total en los factores.
	
	\begin{ejemplo}
		\[
		2.4\times0.3=0.72.
		\]
		Como $24\times3=72$ y hay dos cifras decimales en total,
		el resultado es $0.72$.
	\end{ejemplo}
	
	\section{Porcentajes}
	
	Porcentaje significa ``por cada cien''.
	
	\begin{formula}
		\[
		p\%=\frac{p}{100},\qquad
		p\%\text{ de }N=\frac{p}{100}N.
		\]
	\end{formula}
	
	\begin{ejemplo}
		¿Cuánto es el 15\% de 800?
		\[
		0.15(800)=120.
		\]
	\end{ejemplo}
	
	Para aumentar una cantidad en $p\%$, multiplica por
	$1+\frac{p}{100}$. Para disminuirla, multiplica por
	$1-\frac{p}{100}$.
	
	\section{Razones, proporciones y regla de tres}
	
	Una razón compara cantidades. Una proporción afirma que
	dos razones son iguales:
	\[
	\frac{a}{b}=\frac{c}{d}
	\quad\Longrightarrow\quad ad=bc.
	\]
	
	\begin{ejemplo}
		Si 3 cuadernos cuestan 45 pesos, ¿cuánto cuestan 5 al mismo
		precio unitario?
		\[
		x=\frac{5(45)}{3}=75.
		\]
	\end{ejemplo}
	
	La regla de tres directa sirve cuando ambas cantidades aumentan
	o disminuyen juntas en la misma proporción. Si una aumenta y
	la otra disminuye, revisa si se trata de proporcionalidad inversa.
	
	% ==================================================
	\chapter{Enteros, potencias y raíces}
	
	\section{Números negativos}
	
	En la recta numérica, los números a la derecha son mayores.
	
	\[
	-5<-2<0<3<7.
	\]
	
	Reglas de signos para multiplicar o dividir:
	
	\begin{center}
		\begin{tabular}{c|c}
			Operación de signos & Resultado\\
			\hline
			$+\cdot+$ & $+$\\
			$+\cdot-$ & $-$\\
			$-\cdot+$ & $-$\\
			$-\cdot-$ & $+$
		\end{tabular}
	\end{center}
	
	Sumar un negativo equivale a restar. Restar un negativo
	equivale a sumar.
	
	\section{Potencias}
	
	\[
	a^n=\underbrace{a\cdot a\cdots a}_{n\text{ factores}}.
	\]
	
	\begin{formula}
		Para las bases y exponentes donde las expresiones estén definidas:
		\[
		a^m a^n=a^{m+n},\qquad
		\frac{a^m}{a^n}=a^{m-n}\quad(a\neq0),
		\]
		\[
		(a^m)^n=a^{mn},\qquad
		(ab)^n=a^nb^n,
		\]
		\[
		a^0=1\quad(a\neq0),\qquad
		a^{-n}=\frac{1}{a^n}\quad(a\neq0).
		\]
	\end{formula}
	
	\section{Raíces}
	
	La raíz cuadrada principal de un número no negativo es el número
	no negativo cuyo cuadrado da el original:
	\[
	\sqrt{49}=7,\qquad \sqrt{a^2}=|a|.
	\]
	
	\begin{ejemplo}
		\[
		\sqrt{72}=\sqrt{36\cdot2}=6\sqrt2.
		\]
	\end{ejemplo}
	
	No se puede separar una suma bajo una raíz de esta manera:
	\[
	\sqrt{a+b}\neq\sqrt a+\sqrt b
	\]
	en general.
	
	\section{Notación científica}
	
	Se escribe un número como
	\[
	a\times10^n,\qquad 1\leq |a|<10.
	\]
	
	Ejemplos:
	\[
	450000=4.5\times10^5,\qquad
	0.00032=3.2\times10^{-4}.
	\]
	
	% ==================================================
	\chapter{Álgebra: letras que representan números}
	
	\section{Expresiones algebraicas}
	
	En $3x+5$, la letra $x$ representa un número.
	El número 3 es el coeficiente; 5 es el término constante.
	
	Solo se reducen términos semejantes:
	\[
	3x+5x=8x,\qquad 3x+5\neq8x.
	\]
	
	\section{Propiedad distributiva}
	
	\begin{formula}
		\[
		a(b+c)=ab+ac
		\]
		\[
		a(b-c)=ab-ac
		\]
	\end{formula}
	
	\begin{ejemplo}
		\[
		3(x+4)=3x+12.
		\]
	\end{ejemplo}
	
	\section{Productos notables}
	
	\begin{formula}
		\[
		(a+b)^2=a^2+2ab+b^2
		\]
		\[
		(a-b)^2=a^2-2ab+b^2
		\]
		\[
		(a+b)(a-b)=a^2-b^2
		\]
		\[
		(a+b)^3=a^3+3a^2b+3ab^2+b^3
		\]
	\end{formula}
	
	\begin{ojo}
		$(a+b)^2$ no es $a^2+b^2$. Falta el término $2ab$.
	\end{ojo}
	
	\section{Factorización}
	
	Factorizar es expresar una suma como producto.
	
	\begin{formula}
		\[
		ab+ac=a(b+c)
		\]
		\[
		a^2-b^2=(a-b)(a+b)
		\]
		\[
		x^2+(m+n)x+mn=(x+m)(x+n)
		\]
	\end{formula}
	
	Ejemplo:
	\[
	x^2+5x+6=(x+2)(x+3).
	\]
	Se buscan dos números que sumen 5 y multipliquen 6.
	
	\section{Ecuaciones lineales}
	
	Una ecuación afirma que dos expresiones son iguales.
	Resolverla consiste en encontrar los valores que la hacen cierta.
	
	\begin{ejemplo}
		Resuelve $3x+5=20$.
		\[
		3x=20-5=15,\qquad x=5.
		\]
	\end{ejemplo}
	
	La operación que haces en un lado debes hacerla también
	en el otro. La igualdad no es decoración.
	
	\section{Ecuaciones con paréntesis y fracciones}
	
	Ejemplo:
	\[
	2(x-3)=10
	\]
	\[
	2x-6=10,\qquad 2x=16,\qquad x=8.
	\]
	
	Si aparecen fracciones, puedes multiplicar toda la ecuación
	por el mínimo común múltiplo de los denominadores, siempre
	respetando las restricciones del problema.
	
	\section{Sistemas de ecuaciones}
	
	Un sistema contiene dos o más ecuaciones con las mismas
	incógnitas.
	
	\begin{ejemplo}
		\[
		\begin{cases}
			x+y=7\\
			x-y=1
		\end{cases}
		\]
		
		Sumamos las ecuaciones:
		\[
		2x=8\Longrightarrow x=4.
		\]
		Sustituimos:
		\[
		4+y=7\Longrightarrow y=3.
		\]
	\end{ejemplo}
	
	Métodos frecuentes: sustitución, igualación y reducción.
	
	\section{Ecuaciones cuadráticas}
	
	Forma general:
	\[
	ax^2+bx+c=0,\qquad a\neq0.
	\]
	
	\begin{formula}
		\[
		x=\frac{-b\pm\sqrt{b^2-4ac}}{2a}.
		\]
	\end{formula}
	
	El discriminante es $\Delta=b^2-4ac$:
	\begin{itemize}
		\item $\Delta>0$: dos soluciones reales distintas.
		\item $\Delta=0$: una solución real doble.
		\item $\Delta<0$: no hay soluciones reales.
	\end{itemize}
	
	\begin{ejemplo}
		Resuelve $x^2-5x+6=0$.
		\[
		(x-2)(x-3)=0.
		\]
		Por tanto, $x=2$ o $x=3$.
	\end{ejemplo}
	
	\section{Desigualdades}
	
	Los signos son $<$, $>$, $\leq$ y $\geq$.
	
	Al multiplicar o dividir una desigualdad por un número
	negativo, se invierte el signo.
	
	\begin{ejemplo}
		\[
		-2x>6\Longrightarrow x<-3.
		\]
	\end{ejemplo}
	
	% ==================================================
	\chapter{Funciones, sucesiones y gráficas}
	
	\section{Plano cartesiano}
	
	Un punto se escribe $(x,y)$.
	El eje horizontal es $x$; el vertical es $y$.
	
	Los cuadrantes se numeran comenzando por el superior derecho
	y siguiendo en sentido contrario a las agujas del reloj.
	
	\section{Funciones}
	
	Una función asigna a cada entrada válida exactamente una salida.
	Se escribe $y=f(x)$.
	
	\begin{ejemplo}
		Si $f(x)=2x+1$, entonces
		\[
		f(3)=2(3)+1=7.
		\]
	\end{ejemplo}
	
	\section{Rectas}
	
	La forma pendiente-intersección es
	\[
	y=mx+b,
	\]
	donde $m$ es la pendiente y $b$ el corte con el eje $y$.
	
	Con dos puntos:
	\begin{formula}
		\[
		m=\frac{y_2-y_1}{x_2-x_1},\qquad x_2\neq x_1.
		\]
	\end{formula}
	
	\section{Sucesiones}
	
	Una sucesión es una lista ordenada de números.
	
	Aritmética: se suma siempre la misma diferencia $d$.
	\[
	a_n=a_1+(n-1)d.
	\]
	
	Geométrica: se multiplica siempre por la misma razón $r$.
	\[
	a_n=a_1r^{n-1}.
	\]
	
	Ejemplo aritmético: $3,7,11,15,\ldots$ tiene diferencia 4.
	
	% ==================================================
	\chapter{Geometría plana}
	
	\section{Ángulos}
	
	\begin{itemize}
		\item Agudo: menor que $90^\circ$.
		\item Recto: $90^\circ$.
		\item Obtuso: mayor que $90^\circ$ y menor que $180^\circ$.
		\item Llano: $180^\circ$.
		\item Completo: $360^\circ$.
	\end{itemize}
	
	Los ángulos interiores de un triángulo suman $180^\circ$.
	Los de un cuadrilátero suman $360^\circ$.
	
	\section{Perímetros y áreas}
	
	\begin{center}
		\begin{tabularx}{\textwidth}{lXX}
			\toprule
			Figura & Perímetro & Área\\
			\midrule
			Cuadrado de lado $l$ & $4l$ & $l^2$\\
			Rectángulo & $2(b+h)$ & $bh$\\
			Triángulo & suma de lados & $\frac{bh}{2}$\\
			Paralelogramo & $2(a+b)$ & $bh$\\
			Trapecio & suma de lados & $\frac{(B+b)h}{2}$\\
			Círculo & $2\pi r$ & $\pi r^2$\\
			\bottomrule
		\end{tabularx}
	\end{center}
	
	En las fórmulas, $b$ suele representar base, $h$ altura,
	$r$ radio y $B$ base mayor. Revisa la figura antes de sustituir.
	
	\section{Teorema de Pitágoras}
	
	Para un triángulo rectángulo:
	\[
	a^2+b^2=c^2,
	\]
	donde $c$ es la hipotenusa, el lado opuesto al ángulo recto.
	
	\begin{ejemplo}
		Si los catetos miden 3 y 4:
		\[
		c=\sqrt{3^2+4^2}=\sqrt{25}=5.
		\]
	\end{ejemplo}
	
	No se aplica directamente a cualquier triángulo.
	
	\section{Semejanza y teorema de Tales}
	
	Los triángulos semejantes tienen ángulos correspondientes
	iguales y lados correspondientes proporcionales.
	
	Si las longitudes se escalan por un factor $k$:
	\[
	\text{longitudes}\times k,\qquad
	\text{áreas}\times k^2.
	\]
	
	El teorema de Tales permite establecer proporciones cuando
	rectas paralelas cortan transversales. Identifica primero
	los segmentos correspondientes.
	
	\section{Circunferencia}
	
	\[
	C=2\pi r=\pi d,\qquad A=\pi r^2,\qquad d=2r.
	\]
	
	Una circunferencia es el borde; el círculo es la región interior.
	La terminología importa en los ejercicios.
	
	% ==================================================
	\chapter{Geometría espacial y unidades}
	
	\section{Volúmenes}
	
	\begin{center}
		\begin{tabularx}{\textwidth}{lX}
			\toprule
			Cuerpo & Volumen\\
			\midrule
			Cubo & $V=l^3$\\
			Prisma & $V=A_bh$\\
			Cilindro & $V=\pi r^2h$\\
			Pirámide & $V=\frac13A_bh$\\
			Cono & $V=\frac13\pi r^2h$\\
			Esfera & $V=\frac43\pi r^3$\\
			\bottomrule
		\end{tabularx}
	\end{center}
	
	$A_b$ es el área de la base y $h$ la altura perpendicular.
	
	\section{Conversiones}
	
	En el sistema métrico:
	\[
	1\text{ km}=1000\text{ m},\qquad
	1\text{ m}=100\text{ cm},\qquad
	1\text{ kg}=1000\text{ g}.
	\]
	
	Para áreas, el factor de conversión se eleva al cuadrado;
	para volúmenes, al cubo.
	
	\[
	1\text{ m}^2=10000\text{ cm}^2,\qquad
	1\text{ m}^3=1000000\text{ cm}^3.
	\]
	
	\section{Rapidez, distancia y tiempo}
	
	\begin{formula}
		\[
		v=\frac{d}{t},\qquad d=vt,\qquad t=\frac{d}{v}.
		\]
	\end{formula}
	
	Asegúrate de que las unidades sean compatibles. Kilómetros
	por hora y metros por segundo no son lo mismo.
	
	% ==================================================
	\chapter{Trigonometría}
	
	\section{Razones trigonométricas}
	
	En un triángulo rectángulo, respecto al ángulo $\theta$:
	\begin{formula}
		\[
		\sen\theta=\frac{\text{cateto opuesto}}{\text{hipotenusa}}
		\]
		\[
		\cos\theta=\frac{\text{cateto adyacente}}{\text{hipotenusa}}
		\]
		\[
		\tan\theta=\frac{\text{cateto opuesto}}{\text{cateto adyacente}}
		\]
	\end{formula}
	
	\begin{ejemplo}
		Si el cateto opuesto mide 3 y la hipotenusa 5:
		\[
		\sen\theta=\frac35=0.6.
		\]
	\end{ejemplo}
	
	\section{Identidad fundamental}
	
	\[
	\sen^2\theta+\cos^2\theta=1.
	\]
	
	En una calculadora, revisa si el modo está en grados o radianes.
	Un resultado correcto con la unidad equivocada sigue siendo
	un resultado inútil.
	
	\section{Ley de senos y ley de cosenos}
	
	Para un triángulo con lados $a,b,c$ opuestos a ángulos
	$A,B,C$:
	\[
	\frac{a}{\sen A}=\frac{b}{\sen B}=\frac{c}{\sen C}.
	\]
	
	Ley de cosenos:
	\[
	c^2=a^2+b^2-2ab\cos C.
	\]
	
	Sirven para triángulos que no necesariamente son rectángulos.
	
	% ==================================================
	\chapter{Geometría analítica}
	
	\section{Distancia entre dos puntos}
	
	\[
	d=\sqrt{(x_2-x_1)^2+(y_2-y_1)^2}.
	\]
	
	\section{Punto medio}
	
	\[
	M=\left(\frac{x_1+x_2}{2},\frac{y_1+y_2}{2}\right).
	\]
	
	\section{Ecuación de la circunferencia}
	
	Con centro $(h,k)$ y radio $r$:
	\[
	(x-h)^2+(y-k)^2=r^2.
	\]
	
	\section{Parábola y otras cónicas}
	
	Una forma común de parábola vertical es
	\[
	y=a(x-h)^2+k.
	\]
	
	La elipse y la hipérbola tienen otras ecuaciones canónicas.
	En un examen, identifica primero la forma de la ecuación y
	después qué información pide el problema.
	
	% ==================================================
	\chapter{Estadística descriptiva}
	
	\section{Media, mediana y moda}
	
	Datos: $2,3,3,4,8$.
	
	Media:
	\[
	\bar{x}=\frac{\text{suma de los datos}}{\text{número de datos}}
	=\frac{20}{5}=4.
	\]
	
	Mediana: el dato central al ordenar la lista. Aquí es 3.
	
	Moda: el valor más frecuente. Aquí es 3.
	
	\section{Rango y dispersión}
	
	\[
	\text{Rango}=\text{máximo}-\text{mínimo}.
	\]
	
	La varianza mide dispersión respecto a la media.
	Para una población:
	\[
	\sigma^2=\frac{\sum_{i=1}^n(x_i-\mu)^2}{n}.
	\]
	
	La desviación estándar es la raíz cuadrada de la varianza.
	En muestras, la fórmula usual de varianza muestral divide
	entre $n-1$.
	
	\section{Gráficas}
	
	\begin{itemize}
		\item Barras: comparar categorías.
		\item Histograma: mostrar una distribución numérica agrupada.
		\item Circular: mostrar proporciones de un total.
		\item Línea: observar cambios a lo largo del tiempo.
		\item Dispersión: observar la relación entre dos variables.
	\end{itemize}
	
	Una gráfica puede ser correcta y aun así inducir a error si
	el eje está recortado o las escalas no son comparables.
	
	% ==================================================
	\chapter{Probabilidad y conteo}
	
	\section{Probabilidad clásica}
	
	Cuando todos los resultados son igualmente probables:
	\[
	\Prob(A)=\frac{\text{casos favorables}}{\text{casos posibles}}.
	\]
	
	La probabilidad está entre 0 y 1.
	
	\begin{ejemplo}
		Al lanzar un dado justo, la probabilidad de obtener un 6 es
		\[
		\Prob(6)=\frac16.
		\]
	\end{ejemplo}
	
	\section{Complemento}
	
	\[
	\Prob(A^c)=1-\Prob(A).
	\]
	
	Si la probabilidad de que llueva es 0.3, la de que no llueva
	es 0.7, suponiendo que son los únicos dos resultados.
	
	\section{Principios de conteo}
	
	Si una decisión tiene $m$ opciones y otra tiene $n$ opciones,
	y ambas se pueden combinar libremente, hay $mn$ combinaciones.
	
	Permutaciones de $n$ objetos distintos:
	\[
	n!=n(n-1)(n-2)\cdots2\cdot1.
	\]
	
	Arreglos de $r$ elementos tomados de $n$:
	\[
	P(n,r)=\frac{n!}{(n-r)!}.
	\]
	
	Combinaciones, cuando el orden no importa:
	\[
	\binom nr=\frac{n!}{r!(n-r)!}.
	\]
	
	% ==================================================
	\chapter{Exponenciales y logaritmos}
	
	\section{Función exponencial}
	
	\[
	y=a^x,\qquad a>0,\quad a\neq1.
	\]
	
	Si $a>1$, la función es creciente.
	Si $0<a<1$, es decreciente.
	
	\section{Logaritmos}
	
	El logaritmo responde: ¿a qué exponente hay que elevar
	la base para obtener cierto número?
	
	\[
	\log_a x=y\quad\Longleftrightarrow\quad a^y=x.
	\]
	
	Condiciones: $a>0$, $a\neq1$ y $x>0$.
	
	\begin{formula}
		\[
		\log_a(xy)=\log_a x+\log_a y
		\]
		\[
		\log_a\left(\frac{x}{y}\right)=\log_a x-\log_a y
		\]
		\[
		\log_a(x^r)=r\log_a x
		\]
	\end{formula}
	
	Cambio de base:
	\[
	\log_a x=\frac{\log_b x}{\log_b a}.
	\]
	
	No existe una regla que convierta automáticamente
	$\log(x+y)$ en $\log x+\log y$.
	
	% ==================================================
	\chapter{Límites, derivadas e integrales}
	
	\section{Límites}
	
	Un límite describe el valor al que se aproxima una función
	cuando la variable se acerca a cierto punto.
	
	\[
	\lim_{x\to a}f(x)=L.
	\]
	
	Si la función es continua en $a$, normalmente basta con
	sustituir $x=a$. Si aparece una indeterminación como $0/0$,
	hace falta simplificar o aplicar otro método.
	
	\begin{ejemplo}
		\[
		\lim_{x\to2}(x^2+1)=2^2+1=5.
		\]
	\end{ejemplo}
	
	\section{Derivadas}
	
	La derivada representa una tasa de cambio instantánea.
	Geométricamente, es la pendiente de la recta tangente.
	
	\begin{formula}
		Reglas básicas:
		\[
		\frac{d}{dx}(c)=0,\qquad
		\frac{d}{dx}(x^n)=nx^{n-1}
		\]
		\[
		\frac{d}{dx}(cf)=cf',\qquad
		\frac{d}{dx}(f+g)=f'+g'
		\]
		\[
		(fg)'=f'g+fg'
		\]
		\[
		\left(\frac fg\right)'=
		\frac{f'g-fg'}{g^2},\quad g\neq0.
		\]
	\end{formula}
	
	\begin{ejemplo}
		Si $f(x)=3x^2+2x-4$:
		\[
		f'(x)=6x+2.
		\]
	\end{ejemplo}
	
	\section{Integral}
	
	La integral definida puede representar área con signo
	bajo una curva y acumular cantidades.
	
	\begin{formula}
		\[
		\int x^n\,dx=\frac{x^{n+1}}{n+1}+C,\qquad n\neq-1.
		\]
		\[
		\int_a^b f(x)\,dx=F(b)-F(a),\qquad F'=f.
		\]
	\end{formula}
	
	\begin{ejemplo}
		\[
		\int_0^2 x\,dx
		=\left[\frac{x^2}{2}\right]_0^2
		=2.
		\]
	\end{ejemplo}
	
	Estos temas suelen pertenecer al bachillerato avanzado
	o a cursos introductorios de cálculo; no todos los planes
	de preparatoria los incluyen.
	
	% ==================================================
	\chapter{Errores frecuentes}
	
	\begin{enumerate}
		\item Sumar fracciones sumando también los denominadores.
		\item Olvidar la jerarquía de operaciones.
		\item Cambiar un signo al pasar términos de una ecuación.
		\item Olvidar invertir una desigualdad al dividir entre un negativo.
		\item Confundir radio con diámetro.
		\item Usar área cuando el problema pide perímetro.
		\item Usar grados cuando la calculadora está en radianes.
		\item Confundir una relación proporcional con una relación cualquiera.
		\item Redondear demasiado pronto.
		\item Dar una respuesta sin unidades cuando sí hacen falta.
		\item Dividir entre cero.
		\item Afirmar que $\sqrt{a+b}=\sqrt a+\sqrt b$.
	\end{enumerate}
	
	La comprobación más barata es sustituir el resultado en el
	problema original. Hazla antes de entregar.
	
	% ==================================================
	\chapter{Formulario exprés}
	
	\small
	\begin{multicols}{2}
		
		\textbf{Porcentaje}
		\[
		p\%\text{ de }N=\frac{pN}{100}
		\]
		
		\textbf{Proporción}
		\[
		\frac ab=\frac cd\Rightarrow ad=bc
		\]
		
		\textbf{Pitágoras}
		\[
		a^2+b^2=c^2
		\]
		
		\textbf{Recta}
		\[
		y=mx+b
		\]
		
		\textbf{Pendiente}
		\[
		m=\frac{y_2-y_1}{x_2-x_1}
		\]
		
		\textbf{Cuadrática}
		\[
		x=\frac{-b\pm\sqrt{b^2-4ac}}{2a}
		\]
		
		\textbf{Círculo}
		\[
		A=\pi r^2,\quad C=2\pi r
		\]
		
		\textbf{Triángulo}
		\[
		A=\frac{bh}{2}
		\]
		
		\textbf{Rectángulo}
		\[
		A=bh,\quad P=2b+2h
		\]
		
		\textbf{Prisma}
		\[
		V=A_bh
		\]
		
		\textbf{Cilindro}
		\[
		V=\pi r^2h
		\]
		
		\textbf{Cono}
		\[
		V=\frac13\pi r^2h
		\]
		
		\textbf{Rapidez}
		\[
		v=\frac dt
		\]
		
		\textbf{Media}
		\[
		\bar{x}=\frac{\sum x_i}{n}
		\]
		
		\textbf{Probabilidad clásica}
		\[
		\Prob(A)=\frac{F}{N}
		\]
		
		\textbf{Sucesión aritmética}
		\[
		a_n=a_1+(n-1)d
		\]
		
		\textbf{Sucesión geométrica}
		\[
		a_n=a_1r^{n-1}
		\]
		
		\textbf{Trigonometría}
		\[
		\sen\theta=\frac{CO}{H}
		\]
		\[
		\cos\theta=\frac{CA}{H}
		\]
		\[
		\tan\theta=\frac{CO}{CA}
		\]
		
		\textbf{Derivada de potencia}
		\[
		(x^n)'=nx^{n-1}
		\]
		
		\textbf{Integral de potencia}
		\[
		\int x^n dx=\frac{x^{n+1}}{n+1}+C
		\]
		
	\end{multicols}
	\normalsize
	
	% ==================================================
	\chapter{Mini examen de supervivencia}
	
	Intenta resolver sin mirar las respuestas.
	
	\begin{enumerate}
		\item Calcula $12-2(3+1)$.
		\item Simplifica $\frac{18}{24}$.
		\item Calcula el 20\% de 350.
		\item Resuelve $5x-7=18$.
		\item Factoriza $x^2-9$.
		\item Calcula la hipotenusa con catetos 6 y 8.
		\item Calcula el área de un círculo de radio 3.
		\item Encuentra la media de 4, 6, 8 y 10.
		\item Calcula la pendiente entre $(1,2)$ y $(3,8)$.
		\item Deriva $f(x)=4x^3$.
	\end{enumerate}
	
	\section*{Respuestas}
	\addcontentsline{toc}{section}{Respuestas del mini examen}
	
	\begin{enumerate}
		\item $4$.
		\item $\frac34$.
		\item $70$.
		\item $x=5$.
		\item $(x-3)(x+3)$.
		\item $10$.
		\item $9\pi$ unidades cuadradas.
		\item $7$.
		\item $m=3$.
		\item $f'(x)=12x^2$.
	\end{enumerate}
	
	\vfill
	\begin{center}
		\textbf{Fin.}
		
		Ya tienes las fórmulas. El resto es identificar cuál
		te están pidiendo y no pelearte con los signos.
	\end{center}
	
\end{document}
